% ==============================================================================
% ch:zaver — závěr: reflexe naplnění dílčích cílů, limity, výhled
% Bez nových citací; pouze odkazy na vlastní kapitoly.
% ==============================================================================

Diplomová práce zkoumala, zda slabé pokrytí \aclpins{KS}, nízká odborová organizovanost a~podfinancovaná \acl{APZ} zvyšují zranitelnost českého trhu práce vůči technologickým, mzdovým a~demografickým tlakům. Stanovený hlavní cíl i~navazující dílčí cíle, formulované v~kap.~\ref{par:cile}, byly postupným systematickým postupem plně naplněny. V~teoretické části byl vymezen historický, právní a~pojmový rámec sociálního dialogu; empirická část následně popsala specifika českého trhu práce včetně argumentů svědčících v~jeho prospěch, mezinárodní komparace zasadila \aclacc{ČR} do~kontextu evropských modelů, korelační analýza společně s~modelovými výpočty prověřila provázanost mezi pokrytím \aclpins{KS}, produktivitou práce, délkou pracovní doby a~kupní silou a~provedené případové studie závěrem demonstrovaly, jakým způsobem se stav sociálního dialogu propisuje do~konkrétního obsahu smluv.\par
V~rámci této systematické logiky byly současně zodpovězeny všechny formulované výzkumné otázky, počínaje diagnostikou institucionální síly \aclgen{KV}, přes analýzu kvality konkrétních smluvních výsledků a~identifikaci strukturálních problémů trhu práce, až~po~komplexní hodnocení připravenosti české ekonomiky na~nadcházející technologické, demografické a~odvětvové strukturální změny.\par

Provedená mezinárodní komparace s~\aclins{DK}, \aclins{AT} a~osvojenou praxí v~\aclloc{DE} poukázala na~možnou selektivní přenositelnost jejich funkčních institucionálních mechanismů. Z~dánského modelu flexicurity vyplývá organická vazba mezi tržní flexibilitou a~robustním systémem sociálních jistot. Rakouská praxe ukazuje možnost oddělit vysokou míru pokrytí \aclins{KS} od~samotné odborové organizovanosti, zatímco německá zkušenost zvýrazňuje jak význam odvětvové mzdové koordinace, tak~rizika spojená s~deregulací nízkopříjmové práce. Postkomunistické srovnání v~oblasti střední a~východní Evropy ukázalo, že~historicky oslabené podnikové vyjednávání se bez aktivní legislativní a~politické podpory ze~strany státu nepřeměnilo ve~funkční odvětvovou architekturu. Důležité poznatky přináší konfrontace \aclgen{ČR} s~vývojem v~\aclloc{PL} -- obě ekonomiky silně vázány na subdodavatelské řetězce německého průmyslu, Polsko ale vykazuje rychlejší konvergenci hospodářství v~důsledku silné spotřeby, sycené dynamickým růstem mezd.\par

Formulovaná hlavní hypotéza práce je plně v~souladu s~empirickými zjištěními, a~to v~intencích její verifikace provedené v~kap.~\ref{par:hypoteza_overeni}: \emph{Posílení sociálního dialogu a~kolektivního vyjednávání v~\aclloc{ČR} by přispělo k~udržitelnému hospodářskému rozvoji.} Historicko-institucionální deskripce, mezinárodní komparace, kvantitativní korelační analýza i~hloubkové případové studie shodně demonstrují, že~nízké pokrytí zaměstnanců kolektivními smlouvami je~úzce spojeno s~nedostatečným přenosem produktivitních zisků do~mzdového růstu a~se~zvýšeným rizikem socioekonomické nerovnosti; institucionální posílení vyjednávacího procesu tudíž přímo podporuje strategické osy udržitelného hospodářského rozvoje. Korelace sama o~sobě neprokazuje jednostrannou kauzalitu, nicméně zjištěný směr vykazuje konzistenci s~kvalitativní komparací i~zjištěními případových studií.\par

Z~těchto analytických pilířů bylo v~návrhové části odvozeno pět provázaných systémových inovací. První z~nich, spočívající v~aktivaci mechanismu rozšiřování závaznosti \aclpgen{KSVS} dle \mbox{§\,7} zákona o~\aclloc{KV} zavedením transparentních a~předvídatelných procesů pro rozšiřování závaznosti, bezprostředně cílí na~překonání nízkého pokrytí. Druhý návrh, integrující sociální partnery do~správy a~financování \aclgen{APZ}, rozšiřuje funkční pole sociálního dialogu za~tradiční hranice mzdového vyjednávání směrem k~aktivnímu řízení strukturálních rekvalifikací podle potřeb v~jednotlivých odvětvích. Třetí inovace v~podobě mzdové transparentnosti využívá závazné odvětvové tarify jako klíčový referenční nástroj pro uplatňování zásady stejné odměny za~práci stejné hodnoty. Čtvrtý směr, upravující právní postavení \acpgen{OSVČ} vykonávajících de~facto závislou činnost a~pracovníků digitálních platforem, chrání základnu \aclgen{KV} před erozí. Pátá inovace, zřízení veřejného digitálního registru \aclpgen{KS}, buduje nezbytnou informační infrastrukturu, bez níž nelze zvyšování pokrytí ani mzdovou rovnost kontinuálně vyhodnocovat. Ve~svém souhrnu tyto návrhy neusilují o~mechanickou restituci historického modelu odborového hnutí, nýbrž představují promyšlený pokus rekonfigurovat sociální dialog pro moderní trh práce, v~němž se protínají standardní zaměstnanecký poměr a~samostatné podnikání s~průběžným celoživotním vzděláváním a~technologickou substitucí.\par

%Postup práce vykazuje čtyři hlavní metodologické limity. Zaprvé, korelační analýza operuje s~ekonometrickým panelem členských států \acsgen{EU} v~časově ohraničeném horizontu s~přetrvávajícími mezerami v~datové bázi, čímž nemůže jednoznačně potvrdit sílu a~směr kauzálního působení. Zadruhé, provedené legislativně-modelové výpočty pracují se~statickou strukturou trhu práce a~z~povahy modelu nepostihují dynamické zpětnovazebné reakce na~straně poptávky po~práci či individuální nabídky pracovníků. Zatřetí, vybrané případové studie sice poskytují detailní pohled, nicméně jejich poznatky nelze zobecnit na~celý trh práce. Začtvrté, práce neprovádí analýzu českého daňového mixu, zejména hloubkové rozpadové posouzení reálné míry zdanění a~odvodové zátěže mezi zaměstnanci, zaměstnavateli, \acpins{OSVČ}, kapitálem a~spotřebními daněmi.\par
